\documentclass[11pt,a4paper]{article}
\usepackage[T1]{fontenc}
\usepackage[utf8]{inputenc}
\usepackage{amsmath,amssymb,amsthm}
\usepackage[margin=1in]{geometry}
\usepackage[colorlinks=true,linkcolor=blue,urlcolor=blue]{hyperref}
\theoremstyle{plain}
\newtheorem{theorem}{Theorem}
\newtheorem{lemma}[theorem]{Lemma}
\newtheorem{proposition}[theorem]{Proposition}
\newtheorem{corollary}[theorem]{Corollary}
\theoremstyle{definition}
\newtheorem{remark}[theorem]{Remark}
\newcommand{\Chat}{\widehat{C}}

\title{Carry periodicity forces a large $(x^3+1)$ factor:\\
a proof of Karttunen's conjecture on the binary\\
columns of the Fibonacci numbers}
\author{Adrian Perez Fontelles\\ \small Independent researcher}
\date{23 August 2026}
\begin{document}
\maketitle

\begin{abstract}
Write the Fibonacci numbers one under another in binary. The $n$-th column from the right
is periodic in $k$ with period $3\cdot2^{n}$; reading one full period as the coefficient
vector of a polynomial over $\mathrm{GF}(2)$ gives OEIS A036284. In November 1998
A.~Karttunen conjectured that this polynomial is divisible by $(x^{3}+1)^{2^{n-1}-1}$. We
prove it. The whole content is that the carry into position $n$ depends only on the low
$n$ bits of the two summands and therefore has period $3\cdot2^{n-1}$ --- exactly half the
period of the column itself. Since $x^{3\cdot2^{m}}+1=(x^{3}+1)^{2^{m}}$ in characteristic
2, halving the period is the same as extracting a factor $(x^{3}+1)^{2^{n-1}}$, and the
conjecture drops out of a one-line valuation count. The argument in fact gives the exact
$(x^{2}+x+1)$-adic valuation, not just the lower bound.
\end{abstract}

\section{The sequence and the conjecture}

Let $F(0)=0$, $F(1)=1$, $F(k)=F(k-1)+F(k-2)$. Writing the $F(k)$ in binary, one under
another, the bits in the $n$-th column from the right repeat with period $3\cdot2^{n}$.
OEIS A036284, contributed by Antti Karttunen on 1 November 1998 under the name
\emph{Periodic vertical binary vectors of Fibonacci numbers}, is one full period of that
column read as a binary number:
\[
  a(n)=\sum_{k=0}^{3\cdot2^{n}-1}\Bigl(\bigl\lfloor F(k)/2^{n}\bigr\rfloor\bmod2\Bigr)2^{k},
  \qquad a(0),a(1),a(2),a(3)=6,\;24,\;1440,\;5728448.
\]
Identifying an integer with the $\mathrm{GF}(2)$ polynomial whose coefficient of $x^{k}$ is
bit $k$, the entry carries the comment

\begin{quote}
\textbf{Conjecture:} For $n\ge1$, each term $a(n)$, when considered as a
$\mathrm{GF}(2)[X]$-polynomial, is divisible by the $\mathrm{GF}(2)[X]$-polynomial
$(x^{3}+1)^{\mathrm{A000225}(n-1)}$. \dots\ \textbf{Conjecture 2:} there is also one extra
$(x^{1}+1)$ factor present.
\end{quote}

Here $\mathrm{A000225}(n-1)=2^{n-1}-1$. This note proves the first conjecture. As of the
``Last modified'' line on the live entry (23 August 2026) both statements are still
recorded as unproven conjectures, no proof appears in the entry's links or programs, and a
literature search on binary columns of the Fibonacci numbers turns up nothing.

\section{Notation}

Throughout $n\ge1$ is fixed and
\[
  L=3\cdot2^{n},\qquad q(x)=x^{2}+x+1,\qquad x^{3}+1=(x+1)\,q(x).
\]
For a nonzero $P\in\mathrm{GF}(2)[x]$ and an irreducible $f$, write $v_{f}(P)$ for the
exponent of $f$ in $P$. Two facts are used repeatedly:
\begin{equation}\label{eq:frob}
  x^{3\cdot2^{m}}+1=(x^{3}+1)^{2^{m}}\quad\text{in }\mathrm{GF}(2)[x]\ (m\ge0),
\end{equation}
by the Frobenius endomorphism, and
$v_{q}\bigl((x^{3}+1)^{m}\bigr)=v_{x+1}\bigl((x^{3}+1)^{m}\bigr)=m$.

Let $r_{k}=F(k)\bmod2^{n+1}$, let
\[
  b(k)=\bigl\lfloor r_{k}/2^{n}\bigr\rfloor\in\{0,1\}
\]
be the $n$-th bit, and let $s_{k}=r_{k}\bmod2^{n}$ be the low part, so that
$r_{k}=b(k)2^{n}+s_{k}$. Finally set
\[
  \gamma(k)=\bigl\lfloor(s_{k-1}+s_{k-2})/2^{n}\bigr\rfloor\in\{0,1\},
\]
the carry into position $n$ when $F(k-1)$ and $F(k-2)$ are added. Put
\[
  A(x)=\sum_{k=0}^{L-1}b(k)\,x^{k},
\]
so that $A$ is the polynomial attached to $a(n)$.

\section{Three lemmas}

\begin{lemma}\label{lem:pisano}
For every $m\ge1$ the sequence $\bigl(F(k)\bmod2^{m}\bigr)_{k\ge0}$ is purely periodic with
least period $3\cdot2^{m-1}$.
\end{lemma}

\begin{proof}
Pure periodicity holds for any modulus because the recurrence is invertible:
$F(k-2)=F(k)-F(k-1)$. The value $\pi(2^{m})=3\cdot2^{m-1}$ of the Pisano period at prime
power $2^{m}$ is classical; see Wall~\cite{wall}.
\end{proof}

In particular $\bigl(r_{k}\bigr)_{k}$ has period $L=3\cdot2^{n}$, so $b$, $s$, $\gamma$ are
all well defined as functions on $\mathbb{Z}/L\mathbb{Z}$, and
$r_{k}\equiv r_{k-1}+r_{k-2}\pmod{2^{n+1}}$ holds for every
$k\in\mathbb{Z}/L\mathbb{Z}$ (for $k\ge2$ it is the recurrence; for $k=0,1$ it follows by
periodicity from the instances $k=L$, $L+1$).

\begin{lemma}[carry identity]\label{lem:carry}
For every $k\in\mathbb{Z}/L\mathbb{Z}$,
\[
  b(k)+b(k-1)+b(k-2)\equiv\gamma(k)\pmod2.
\]
\end{lemma}

\begin{proof}
Split both summands into high bit and low part:
\[
  r_{k-1}+r_{k-2}=\bigl(b(k-1)+b(k-2)\bigr)2^{n}+\bigl(s_{k-1}+s_{k-2}\bigr)
  =\bigl(b(k-1)+b(k-2)+\gamma(k)\bigr)2^{n}+\bigl(s_{k-1}+s_{k-2}\bigr)\bmod2^{n},
\]
using $s_{k-1}+s_{k-2}=\gamma(k)2^{n}+\bigl((s_{k-1}+s_{k-2})\bmod2^{n}\bigr)$, which is
legitimate because $s_{k-1}+s_{k-2}<2^{n+1}$. Reducing modulo $2^{n+1}$ and comparing with
$r_{k}=b(k)2^{n}+s_{k}$, where $s_{k}=(s_{k-1}+s_{k-2})\bmod2^{n}<2^{n}$, forces
$b(k)\equiv b(k-1)+b(k-2)+\gamma(k)\pmod2$.
\end{proof}

\begin{lemma}[the carry has half the period]\label{lem:half}
$\gamma(k+3\cdot2^{n-1})=\gamma(k)$ for every $k$.
\end{lemma}

\begin{proof}
By definition $\gamma(k)$ is a function of $s_{k-1}$ and $s_{k-2}$ alone, and
$s_{j}=F(j)\bmod2^{n}$. By Lemma~\ref{lem:pisano} with $m=n$ the sequence $(s_{j})_{j}$ has
period $3\cdot2^{n-1}$.
\end{proof}

This is the whole point: the column $b$ has period $L=3\cdot2^{n}$, while the carry
$\gamma$ that drives it has period only $L/2$.

\section{The theorem}

\begin{theorem}\label{thm:main}
Let $n\ge1$ and let $\Chat(x)=\sum_{k=0}^{L/2-1}\gamma(k)x^{k}$. Then in
$\mathrm{GF}(2)[x]$
\begin{equation}\label{eq:main}
  q(x)\,A(x)=\Chat(x)\,(x^{3}+1)^{2^{n-1}}+M(x)\,(x^{3}+1)^{2^{n}}
\end{equation}
for some $M\in\mathrm{GF}(2)[x]$. Consequently
\[
  v_{q}(A)\ge2^{n-1}-1,\qquad v_{x+1}(A)\ge2^{n-1},
\]
and therefore $(x^{3}+1)^{2^{n-1}-1}$ divides $A$.
\end{theorem}

\begin{proof}
Multiply the identity of Lemma~\ref{lem:carry} by $x^{k}$ and sum over
$k\in\mathbb{Z}/L\mathbb{Z}$. On the left, reading indices modulo $L$ turns the shift
$k\mapsto k-1$ into multiplication by $x$ modulo $x^{L}+1$, so the left side becomes
$(1+x+x^{2})A(x)=q(x)A(x)$ modulo $x^{L}+1$. On the right we get
$\Gamma(x)=\sum_{k=0}^{L-1}\gamma(k)x^{k}$. Hence
\[
  q(x)A(x)\equiv\Gamma(x)\pmod{x^{L}+1}.
\]
By Lemma~\ref{lem:half} the coefficient list of $\Gamma$ is the concatenation of two copies
of that of $\Chat$, that is
\[
  \Gamma(x)=\Chat(x)\bigl(1+x^{L/2}\bigr)=\Chat(x)\bigl(x^{3\cdot2^{n-1}}+1\bigr)
  =\Chat(x)\,(x^{3}+1)^{2^{n-1}},
\]
the last step by \eqref{eq:frob}. Since $x^{L}+1=(x^{3}+1)^{2^{n}}$, again by
\eqref{eq:frob}, this is exactly \eqref{eq:main}.

Now take valuations in \eqref{eq:main}. Both terms on the right are divisible by
$(x^{3}+1)^{2^{n-1}}$, hence
\[
  v_{q}\bigl(qA\bigr)\ge2^{n-1},\qquad v_{x+1}\bigl(qA\bigr)\ge2^{n-1}.
\]
Since $v_{q}(q)=1$ and $v_{x+1}(q)=0$, this reads $v_{q}(A)\ge2^{n-1}-1$ and
$v_{x+1}(A)\ge2^{n-1}$. Finally
$v_{x^{3}+1}(A)=\min\bigl(v_{x+1}(A),v_{q}(A)\bigr)\ge2^{n-1}-1$.
\end{proof}

\begin{corollary}
For every $n\ge1$, $(x^{3}+1)^{2^{n-1}-1}$ divides $a(n)$ in $\mathrm{GF}(2)[X]$. This is
Karttunen's conjecture.
\end{corollary}

\begin{remark}[the exact valuation]
Equation \eqref{eq:main} gives more than the bound. If $v_{q}(\Chat)<2^{n-1}$ then the two
right-hand terms have different $q$-valuations and
\[
  v_{q}(A)=2^{n-1}-1+v_{q}(\Chat),
\]
and likewise $v_{x+1}(A)=2^{n-1}+v_{x+1}(\Chat)$ when $v_{x+1}(\Chat)<2^{n-1}$.
Computation (Section~5) gives $v_{q}(\Chat)=0$ for $1\le n\le12$, so $v_{q}(A)=2^{n-1}-1$
exactly there --- the conjectured bound is sharp on the $q$-side. On the $(x+1)$-side
$v_{x+1}(\Chat)=0$ for $n=1$ and $4\le n\le12$, but $v_{x+1}(\Chat)=1$ at $n=2$ and 3 at
$n=3$, which is why $v_{x+1}(A)$ equals 3 and 7 there instead of 2 and 4. I have not proved
that $v_{q}(\Chat)=0$ for all $n$; that statement is equivalent to $\Chat(\omega)\ne0$ for
$\omega$ a primitive cube root of unity in $\mathrm{GF}(4)$, that is to a condition on the
distribution modulo 3 of the positions of the carries. It is an open question here.
\end{remark}

\begin{remark}[what the proof really uses]
Nothing about the Fibonacci recurrence is used except that the modulus-$2^{m}$ period
doubles when $m$ increases by one. The same argument gives, verbatim, the analogous
statement for any integer sequence satisfying $u(k)=u(k-1)+u(k-2)$ whose period modulo
$2^{m}$ is $\pi\cdot2^{m-1}$ for a fixed odd $\pi$: the column polynomial is divisible by
$(x^{\pi}+1)^{2^{n-1}-1}$. The proof is elementary and the natural destination is an OEIS
comment, not a journal.
\end{remark}

\section{Verification}

A verification script, written separately from this note, checks every assertion made here
and prints every number quoted. It (i) recomputes $a(0),\dots,a(3)$ from the defining sum
and matches the published values $6,24,1440,5728448$ exactly; (ii) checks
Lemma~\ref{lem:carry} at all $L$ residues for $n=0,\dots,16$, that is up to $L=196608$,
with no failure; (iii) checks Lemma~\ref{lem:half} over the same range, with no failure;
(iv) verifies the identity \eqref{eq:main} by explicit polynomial division for
$n=1,\dots,12$; (v) computes the valuations, obtaining $v_{q}(A)=2^{n-1}-1$ and
$v_{x+1}(A)=2^{n-1}$ for $n=1$ and $4\le n\le12$, and the two exceptional values
$v_{x+1}(A)=3$ at $n=2$ and 7 at $n=3$; and (vi) divides $a(n)$ by $(x^{3}+1)^{2^{n-1}-1}$
directly for $n=1,\dots,12$, obtaining remainder zero in every case.

\begin{thebibliography}{9}
\bibitem{wall} D.~D. Wall, \emph{Fibonacci series modulo $m$}, Amer. Math. Monthly 67
(1960), 525--532.
\bibitem{ln} R.~Lidl and H.~Niederreiter, \emph{Finite Fields}, 2nd ed., Cambridge
University Press, 1997.
\bibitem{oeis} N.~J.~A. Sloane et al., \emph{The On-Line Encyclopedia of Integer
Sequences}, \url{https://oeis.org/A036284}, entry A036284, consulted 23 August 2026.
\end{thebibliography}
\end{document}
